\section{九、项目落地——面向一线的可推广智能监测方案}

本项目已在黑龙江松嫩平原开展试点应用，形成了成熟的落地模式，面向黑土地保护一线，提供“装得起、用得上、测得准”的智能监测工具：“装得起”指硬件门槛足够低，“用得上”指基层人员无需遥感专业背景即可独立操作，“测得准”指真实业务场景下识别结果稳定可靠，三者共同构成项目从技术成果走向一线生产力的闭环。

\subsection{1.试点区域与成效}

项目以黑龙江松嫩平原为试点区域，实现了耕地变化的“周级”动态监测，能够精准识别撂荒、非法占用、侵蚀沟扩展等关键黑土地退化类型，为基层土地管理提供了及时、可靠的决策依据。

松嫩平原地处东北黑土区核心地带，耕地连片规整，既便于卫星影像解译结果与地面核查逐块比对，也覆盖了较为集中的黑土地退化问题，能够使算法在真实的下垫面条件下得到充分检验。在监测频次上，项目改变了以往依赖人工实地巡查、成果按年度更新的作业方式，转为“周级”动态监测，所针对的是两类不同的变化规律：非法占用等行为具有突发性和隐蔽性，一旦错过发现窗口，后续取证与整改成本将显著上升；撂荒与侵蚀沟扩展变化相对缓慢，需要长时间序列的连续观测才能判断其趋势。项目采用的改进 BIT 双时相 Transformer 模型以语义分词器提取语义 tokens，在双时相影像之间建立长程语义对应关系，可对不同尺度、不同形态的变化类型形成统一而敏感的响应，使成果能以较短周期送达基层，为日常巡查与整改跟踪提供靶向指引。

松嫩平原在东北黑土区中具有突出的代表性：黑土成土过程典型、耕层有机质本底较高，同时又是垦殖强度大、人为扰动频繁的地带，撂荒、非法占用与侵蚀沟扩展等退化类型在同一区域内并存，构成了一组天然的“问题样本集”。试点选在此处，便于在相同气候与耕作制度背景下比较各类变化的影像特征差异，避免因跨区域下垫面条件悬殊而干扰对算法性能的判断；区域内耕地权属与地块边界资料相对完整，解译图斑可以直接落到具体地块上，使精度评估能与地面核查结果逐块对应，结论更具说服力。

由年度更新转向周级监测，改变的不只是频次，更是监测在管理链条中的位置。年度成果通常只能在事后给出变化结论，属于“回头看”的统计性描述；周级成果则使变化在发生过程中即被发现，管理者可以在整改窗口尚未关闭时介入，把监测结果直接转化为巡查任务与处置指令。以非法占用为例，此类行为往往在数周内完成并迅速改变地表形态，若按年度复核，待影像比对确认时现场状态可能已难以恢复；周级复核则可在变化初期锁定图斑位置与范围，为现场取证保留更完整的时间证据链。对撂荒与侵蚀沟扩展而言，单周变化幅度往往很小，连续观测能够抑制单期影像噪声带来的误判，使趋势判断更多依赖稳定的时间信号。更高的频次也对算力与存储提出更高要求，这正是系统坚持轻量化、高吞吐设计的现实动因。

\begin{figure}[htbp]
    \centering
    \fbox{\parbox[c][4.5cm][c]{0.8\textwidth}{\centering 【图片空位：image25.png】试点区域技术路线与业务流程示意}}
    \caption{简单的技术路线}
    \label{fig:9}
\end{figure}

\subsection{2.部署方式与适配性}

系统采用前后端轻量化架构，普通电脑即可本地化部署运行，无需专业服务器支持；同时兼容GF-2卫星影像与无人机航拍影像，适配多源数据采集场景，可满足大范围卫星批量监测与重点区域无人机现场核查的双重需求。

在架构上，后端基于 FastAPI 承担影像接收、任务调度、模型推理与结果组织，前端可视化 Web 系统提供影像上传、变化检测、结果可视化、AI 解读与一键报告等功能模块，二者通过标准接口解耦，使算法模型迭代不必牵动前端交互逻辑。轻量化带来的直接价值是部署门槛下降：区县一级的农业技术推广或土地管理机构利用现有办公计算机即可搭建运行，无需专用服务器集群；影像与解译成果在本地完成处理与存储，也减少了敏感地理信息的外链流转。在多源适配方面，GF-2亚米级光学影像覆盖范围广、重访周期稳定，适合大范围批量普查；无人机航拍影像分辨率更高、观测角度灵活，适合对疑点图斑开展现场核查，系统以统一规范组织两类影像，使“卫星普查—无人机核查”能在同一平台内闭环完成。

“卫星普查—无人机核查”两级作业模式的配合，关键在于两类数据在统一空间基准与统一处理规范下衔接。卫星端承担面状覆盖，按周对试点区耕地开展全域比对，输出疑点图斑清单，明确其位置、范围与疑似变化类型；无人机端按清单开展点状核查，对交通不便、地块零碎或影像受云影遮挡的区域低空补测，获取分辨率更高的正射影像与局部细节。两级成果在同一平台内叠加显示，核查结论可以直接回填到卫星图斑上，形成“发现—核验—确认”的固定流程，避免两套数据各自成图、结论难以对应。这种分工在成本上同样合理：面状普查交由重访周期稳定的卫星完成，而受空域审批与天气条件制约的低空作业只投放在真正需要的疑点区域。

\subsection{3.一线友好的使用体验}

系统设计充分考虑基层用户需求，无需遥感专业知识，基层农技人员即可独立操作；性能上，系统采用轻量化模型与批量化推理策略，在保证精度的同时实现了高效率处理，大幅降低了一线监测的人力与时间成本。

从使用流程看，基层人员只需完成影像上传，系统即自动完成变化检测计算、结果渲染与报告生成：变化图斑的位置、范围与类别以可视化图层叠加呈现，无需解读复杂的指数图或波段组合；AI 解读模块以自然语言说明图斑的变化性质与可能成因；一键报告功能将检测结论与关键截图整合为规范化成果文档，使遥感解译的专业门槛被转移至系统内部。在性能上，轻量化模型使交互式操作几乎感受不到等待，批量化推理能力则保证在周级监测频次下足以支撑较大范围区域的批量作业，不会因算力不足而被迫降低监测频次或压缩覆盖范围。

\section{十、项目核心知识产权成果}

本项目在技术研发过程中，形成了以软件著作权与实用新型专利为核心的完整知识产权保护体系，为技术方案的原创性、先进性与可落地性提供了坚实支撑，具体成果如下：

\subsection{1.计算机软件著作权（3项）}

团队成员围绕黑土地遥感监测、耕地变化识别、数据处理与智能分析等核心业务模块，共完成3项计算机软件著作权登记，覆盖了从影像预处理、算法运算到成果可视化的全流程软件系统，充分证明项目软件系统的原创性与独立开发能力，为后续平台迭代与推广应用筑牢了法律基础。

\begin{itemize}
    \item 影像预处理模块：完成辐射校正、几何配准与质量筛查，为算法运算提供规范、可比对的双时相数据基础。
    \item 算法运算模块：承载改进 BIT 双时相 Transformer 模型推理，以 ResNet18为 CNN 骨干网络提取特征，通过语义分词器抽取语义 tokens。
    \item 成果可视化模块：提供影像上传、变化检测、结果可视化、AI 解读与一键报告等交互功能，将算法输出转化为业务决策可用的直观成果。
\end{itemize}

三项登记的完成明确了成果的权利归属，也为技术方案的推广复制保留了稳定的权利基础。

团队成员的计算机软件著作权共计3项，上述内容即为其全部登记成果，不存在其他未列明的登记项。三项登记沿“采集—解译—应用”链条依次落位：影像预处理模块把住数据入口的质量关，解决多源影像在辐射与几何上的可比性问题；算法运算模块承载改进 BIT 双时相 Transformer 的核心推理逻辑；成果可视化模块面向最终使用者，把模型输出转化为可理解、可汇报的业务成果。三者首尾相接、彼此支撑，任一环节缺失都会削弱整条数据链的可用性。

\subsection{2.实用新型专利（1项，已授权）}

项目核心技术方案《一种用于图像采集的遥感无人机》已获国家知识产权局实用新型专利授权，专利号 ZL 2025 2 2047856.0，授权公告号 CN 224511490 U，专利申请日为2025年9月24日，授权公告日为2026年7月17日，专利权人为黑龙江八一农垦大学。该专利聚焦于遥感数据采集与处理的关键创新点，进一步强化了项目在算法与技术实现层面的核心竞争力，为技术成果转化与市场应用构建了知识产权壁垒。

该专利针对的是数据链条最前端的采集环节：在“天空地一体化”体系中，无人机承担高分辨率、高灵活性地获取重点区域影像的职责，其采集质量直接决定后续解译成果的上限。将关键采集方案纳入专利布局，使项目实现了对“采集—解译—应用”全链条中数据入口环节的知识产权覆盖，与算法层面的软件著作权互补，构成从硬件采集方法到软件解译系统的立体保护。

软件著作权与实用新型专利的组合，使保护范围覆盖了从硬件采集方法到软件解译系统的完整链条：专利锁定数据入口环节，保护采集装置与作业方法层面的原创贡献；软件著作权锁定数据处理与成果输出环节，保护代码实现与系统功能层面的原创表达。二者在权利类型上互补，即便他人在硬件上采用不同实现路径，只要复用了本系统的解译与可视化功能，仍会落入软件著作权的保护范围。这种立体布局为后续技术成果转化、平台推广与对外合作划定了可界定的权利边界。

\begin{figure}[htbp]
    \centering
    \fbox{\parbox[c][4.5cm][c]{0.8\textwidth}{\centering 【图片空位：image26.png】实用新型专利证书}}
    \caption{实用新型专利证书（专利号 ZL 2025 2 2047856.0）}
    \label{fig:10}
\end{figure}

\begin{figure}[htbp]
    \centering
    \fbox{\parbox[c][4.5cm][c]{0.8\textwidth}{\centering 【图片空位：image27.png】计算机软件著作权登记证书（一）}}
    \caption{计算机软件著作权登记证书（一）}
    \label{fig:11}
\end{figure}

\begin{figure}[htbp]
    \centering
    \fbox{\parbox[c][4.5cm][c]{0.8\textwidth}{\centering 【图片空位：image28.png】计算机软件著作权登记证书（二）}}
    \caption{计算机软件著作权登记证书（二）}
    \label{fig:12}
\end{figure}
